\begin{textsample}{POS Dim 1 – 2018 – Score 17.00 – 2018/t000534.txt}  \label{ex:f1_pos_068}
All life, every living thing ever, has been built according to \textbf{the} information in DNA.
What does that mean?
Well, it means that just as \textbf{the} English language is made up of alphabetic letters that, when combined into words, allow me to tell you \textbf{the} story I ' m going to tell you today, DNA is made up of genetic letters that, when combined into genes, allow cells to produce proteins, strings of amino acids that fold up into complex structures that perform \textbf{the} functions that allow a cell to do what it does, to tell its stories. \textbf{The} English alphabet has 26 letters, and \textbf{the} genetic alphabet has four.
They ' re pretty \textbf{famous}.
Maybe you ' ve heard of them.
They are often just referred to as G, C, A and T.
But it ' s remarkable that all \textbf{the} diversity of life is \textbf{the} result of four genetic letters.
Imagine what it would be like if \textbf{the} English alphabet had four letters.
What sort of stories would you be able to tell?
What if \textbf{the} genetic alphabet had more letters?
Would life with more letters be able to tell different stories, maybe even more interesting ones?
In 1999, my lab at \textbf{the} Scripps Research Institute in La Jolla, California started working on this question with \textbf{the} goal of creating living organisms with DNA made up of a six-letter genetic alphabet, \textbf{the} four natural letters plus two additional new man-made letters.
Such an organism would be \textbf{the} first radically altered form of life ever created.
It would be a semisynthetic form of life that stores more information than life ever has before.
It would be able to make new proteins, proteins built from more than \textbf{the} 20 normal amino acids that are usually used to build proteins.
What sort of stories could that life tell?
With \textbf{the} power of synthetic \textbf{chemistry} and molecular \textbf{biology} and just under 20 years of work, we created bacteria with six-letter DNA.
Let me tell you how we did it.
All you have to remember from your high school \textbf{biology} is that \textbf{the} four natural letters pair together to form two base pairs.
G pairs with C and A pairs with T, so to create our new letters, we synthesized hundreds of new candidates, new candidate letters, and examined their abilities to selectively pair with each other.
And after about 15 years of work, we found two that paired together really well, at least in a test tube.
They have complicated names, but let ' s just call them X and Y. \textbf{The} next thing we needed to do was find a way to get X and Y into cells, and eventually we found that a protein that does something similar in algae worked in our bacteria.
So \textbf{the} final thing that we needed to do was to show that with X and Y provided, cells could grow and divide and hold on to X and Y in their DNA.
Everything we had done up to then took longer than I had hoped — I am actually a really impatient person — but this, \textbf{the} most important step, worked faster than I dreamed, basically immediately.
On a weekend in 2014, a graduate student in my lab grew bacteria with six- letter DNA.
Let me take \textbf{the} opportunity to introduce you to them right now.
This is an actual picture of them.
These are \textbf{the} first semisynthetic organisms.
So bacteria with six-letter DNA, that ' s really cool, right?
Well, maybe some of you are still wondering why.
So let me tell you a little bit more about some of our motivations, both conceptual and practical.
Conceptually, people have thought about life, what it is, what makes it different from things that are not alive, since people have had thoughts.
Many have interpreted life as being perfect, and this was taken as evidence of a creator.
Living things are different because a god breathed life into them.
Others have sought a more scientific explanation, but I think it ' s fair to say that they still consider \textbf{the} molecules of life to be special.
I mean, \textbf{evolution} has been optimizing them for billions of years, right?
Whatever perspective you take, it would seem pretty impossible for chemists to come in and build new parts that function within and alongside \textbf{the} natural molecules of life without somehow really screwing everything up.
But just how perfectly created or \textbf{evolved} are we?
Just how special are \textbf{the} molecules of life?
These questions have been impossible to even ask, because we ' ve had nothing to compare life to.
Now for \textbf{the} first time, our work suggests that maybe \textbf{the} molecules of life aren ' t that special.
Maybe life as we know it isn ' t \textbf{the} only way it could be.
Maybe we ' re not \textbf{the} only solution, maybe not even \textbf{the} best solution, just a solution.
These questions address fundamental issues about life, but maybe they seem a little esoteric.
So what about practical motivations?
Well, we want to explore what sort of new stories life with an expanded vocabulary could tell, and remember, stories here are \textbf{the} proteins that a cell produces and \textbf{the} functions they have.
So what sort of new proteins with new types of functions could our semisynthetic organisms make and maybe even use?
Well, we have a couple of things in mind. \textbf{The} first is to get \textbf{the} cells to make proteins for us, for our use.
Proteins are being used today for an increasingly broad range of different applications, from materials that protect soldiers from injury to \textbf{devices} that \textbf{detect} dangerous compounds, but at least to me, \textbf{the} most exciting application is protein drugs.
Despite being relatively new, protein drugs have already revolutionized medicine, and, for example, insulin is a protein.
You ' ve probably heard of it, and it ' s manufactured as a drug that has completely changed how we treat diabetes.
But \textbf{the} problem is that proteins are really hard to make and \textbf{the} only practical way to get them is to get cells to make them for you.
So of course, with natural cells, you can only get them to make proteins with \textbf{the} natural amino acids, and so \textbf{the} properties those proteins can have, \textbf{the} applications they could be developed for, must be limited by \textbf{the} nature of those amino acids that \textbf{the} protein ' s built from.
So here they are, \textbf{the} 20 normal amino acids that are strung together to make a protein, and I think you can see, they ' re not that different-looking.
They don ' t bring that many different functions.
They don ' t make that many different functions available.
Compare that with \textbf{the} small molecules that synthetic chemists make as drugs.
Now, they ' re much simpler than proteins, but they ' re routinely built from a much broader range of diverse things.
Don ' t worry about \textbf{the} molecular details, but I think you can see how different they are.
And in fact, it ' s their differences that make them great drugs to treat different diseases.
So it ' s really provocative to wonder what sort of new protein drugs you could develop if you could build proteins from more diverse things.
So can we get our semisynthetic organism to make proteins that include new and different amino acids, maybe amino acids selected to confer \textbf{the} protein with some desired property or function?
For example, many proteins just aren ' t stable when you inject them into people.
They are rapidly degraded or eliminated, and this stops them from being drugs.
What if we could make proteins with new amino acids with things attached to them that protect them from their environment, that protect them from being degraded or eliminated, so that they could be better drugs?
Could we make proteins with little fingers attached that specifically grab on to other molecules?
Many small molecules failed during development as drugs because they just weren ' t specific enough to find their target in \textbf{the} complex environment of \textbf{the} human body.
So could we take those molecules and make them parts of new amino acids that, when incorporated into a protein, are guided by that protein to their target?
I started a biotech company called Synthorx.
Synthorx stands for synthetic organism with an X added at \textbf{the} end because that ' s what you do with biotech companies.
Synthorx is working closely with my lab, and they ' re interested in a protein that recognizes a certain receptor on \textbf{the} \textbf{surface} of human cells.
But \textbf{the} problem is that it also recognizes another receptor on \textbf{the} \textbf{surface} of those same cells, and that makes it toxic.
So could we produce a variant of that protein where \textbf{the} part that interacts with that second bad receptor is shielded, blocked by something like a big umbrella so that \textbf{the} protein only interacts with that first good receptor?
Doing that would be really difficult or impossible to do with \textbf{the} normal amino acids, but not with amino acids that are specifically designed for that purpose.
So getting our semisynthetic cells to act as little factories to produce better protein drugs isn ' t \textbf{the} only potentially really interesting application, because remember, it ' s \textbf{the} proteins that allow cells to do what they do.
So if we have cells that make new proteins with new functions, could we get them to do things that natural cells can ' t do?
For example, could we develop semisynthetic organisms that when injected into a person, seek out cancer cells and only when they find them, secrete a toxic protein that kills them?
Could we create bacteria that eat different kinds of \textbf{oil}, maybe to clean up an \textbf{oil} \textbf{spill}?
These are just a couple of \textbf{the} types of stories that we ' re going to see if life with an expanded vocabulary can tell.
So, sounds great, right?
Injecting semisynthetic organisms into people, dumping millions and millions of gallons of our bacteria into \textbf{the} ocean or out on your favorite beach?
Oh, wait a minute, actually it sounds really scary.
This dinosaur is really scary.
But here ' s \textbf{the} catch : our semisynthetic organisms in order to survive, need to be fed \textbf{the} chemical precursors of X and Y.
X and Y are completely different than anything that exists in nature.
Cells just don ' t have them or \textbf{the} ability to make them.
So when we prepare them, when we grow them up in \textbf{the} controlled environment of \textbf{the} lab, we can feed them lots of \textbf{the} unnatural food.
Then, when we deploy them in a person or out on a beach where they no longer have access that special food, they can grow for a little bit, they can survive for a little, maybe just long enough to perform some intended function, but then they start to run out of \textbf{the} food.
They start to starve.
They starve to death and they just \textbf{disappear}.
So not only could we get life to tell new stories, we get to tell life when and where to tell those stories.
At \textbf{the} beginning of this talk I told you that we reported in 2014 \textbf{the} creation of semisynthetic organisms that store more information, X and Y, in their DNA.
But all \textbf{the} motivations that we just talked about require cells to use X and Y to make proteins, so we started working on that.
Within a couple years, we showed that \textbf{the} cells could take DNA with X and Y and \textbf{copy} it into RNA, \textbf{the} working \textbf{copy} of DNA.
And late last year, we showed that they could then use X and Y to make proteins.
Here they are, \textbf{the} stars of \textbf{the} show, \textbf{the} first fully- functional semisynthetic organisms.
These cells are green because they ' re making a protein that glows green.
It ' s a pretty \textbf{famous} protein, actually, from jellyfish that a lot of people use in its natural form because it ' s easy to see that you made it.
But within every one of these proteins, there ' s a new amino acid that natural life can ' t build proteins with.
Every living cell, every living cell ever, has made every one of its proteins using a four-letter genetic alphabet.
These cells are living and growing and making protein with a six-letter alphabet.
These are a new form of life.
This is a semisynthetic form of life.
So what about \textbf{the} future?
My lab is already working on expanding \textbf{the} genetic alphabet of other cells, including human cells, and we ' re getting ready to start working on more complex organisms.
Think semisynthetic worms. \textbf{The} last thing I want to say to you, \textbf{the} most important thing that I want to say to you, is that \textbf{the} time of semisynthetic life is here.
Thank you.
Chris Anderson : I mean, Floyd, this is so remarkable.
I just wanted to ask you, what are \textbf{the} implications of your work for how we should think about \textbf{the} possibilities for life, like, in \textbf{the} \textbf{universe}, elsewhere?
It just seems like so much of life, or so much of our assumptions are based on \textbf{the} fact that of course, it ' s got to be DNA, but is \textbf{the} possibility space of self-replicating molecules much bigger than DNA, even just DNA with six letters?
Floyd Romesberg : Absolutely, I think that ' s right, and I think what our work has shown, as I mentioned, is that there ' s been always this prejudice that sort of we ' re perfect, we ' re optimal, God created us this way, \textbf{evolution} perfected us this way.
We ' ve made molecules that work right alongside \textbf{the} natural ones, and I think that suggests that any molecules that obey \textbf{the} fundamental laws of \textbf{chemistry} and physics and you can optimize them could do \textbf{the} things that \textbf{the} natural molecules of life do.
There ' s nothing magic there.
And I think that it suggests that life could \textbf{evolve} many different ways, maybe similar to us with other types of DNA, maybe things without DNA at all.
CA : I mean, in your mind, how big might that possibility space be?
Do we even know?
Are most things going to look something like a DNA molecule, or something radically different that can still self-reproduce and potentially create living organisms?
FR : My personal opinion is that if we found new life, we might not even recognize it.
CA : So this obsession with \textbf{the} \textbf{search} for Goldilocks \textbf{planets} in exactly \textbf{the} right place with water and whatever, that ' s a very parochial assumption, perhaps.
FR : Well, if you want to find someone you can talk to, then maybe not, but I think that if you ' re just looking for any form of life, I think that ' s right, I think that you ' re looking for life under \textbf{the} light post.
CA : Thank you for boggling all our minds.
Thank so much, Floyd.

% matched lemmas: biology, chemistry, copy, detect, device, disappear, evolution, evolve, famous, oil, planet, search, spill, surface, the, universe
\end{textsample}
